% GENERATED FILE. Edit canonical lessons, then run npm run generate:books.
% canonical-source-hash: fnv1a64:a0984ec8913b8c7a
% canonical-lessons: MW-C115-jamin, MW-C115-jaydad, MW-C115-puji, MW-C115-kamai, MW-C115-naukri

\chapter{Land, Property, Capital, Earnings, Job}
\label{ch:mw-land-property-capital-earnings-job}
\hlchaptermodality{115}

\begin{chapteropening}
\textbf{By the end of this chapter:} \emph{I can say land, property, capital, earnings and a job.}

Complete the last lesson of chapter 115: I can say land, property, capital, earnings and a job.
\end{chapteropening}

\section[\texorpdfstring{\mw{जमीन} (jamīn)}{जमीन (jamīn)}]{\mw{जमीन} (jamīn) — land}

\label{lesson:MW-C115-jamin}



\begin{quote}
Before the new one: say the Marwadi for courage, then the Marwadi for an owner.
\end{quote}

\subsection*{You'll want to know: \mw{जमीन}}
\textbf{\mw{जमीन}} — \emph{jamīn} — \textquotedblleft{}land\textquotedblright{}.

\begin{cousinweb}
\textbf{\mw{जमीन}} (\emph{jamīn}) is a loan, usually traced to Persian \emph{zamīn}, \textquotedblleft{}earth, ground\textquotedblright{}. Marwadi says the \emph{z} as \emph{j}.
\end{cousinweb}

\begin{grammarlens}[title={What you've built}]
One more word for this chapter.
\end{grammarlens}

\subsection*{Your turn}
\begin{itemize}
\raggedright
  \item \emph{Say it:} \emph{jamīn}
  \item \emph{Say it:} \emph{jamīn}, once more
  \item \emph{Say it:} say \emph{jamīn}
  \item \emph{Recall:} say the Marwadi for courage, then the Marwadi for an owner, then say \emph{jamīn} again
\end{itemize}

\begin{tcolorbox}[breakable,colback=teal!4,colframe=teal!35!black,title={Before you move on}]
What does \mw{जमीन} mean? (\textquotedblleft{}Land\textquotedblright{}.) Say it once more.
\end{tcolorbox}
\section[\texorpdfstring{\mw{जायदाद} (jāydād)}{जायदाद (jāydād)}]{\mw{जायदाद} (jāydād) — property}

\label{lesson:MW-C115-jaydad}



\begin{quote}
Before the new one: say the Marwadi for an owner, then the Marwadi for land.
\end{quote}

\subsection*{You'll want to know: \mw{जायदाद}}
\textbf{\mw{जायदाद}} — \emph{jāydād} — \textquotedblleft{}property\textquotedblright{}.

\begin{cousinweb}
\textbf{\mw{जायदाद}} (\emph{jāydād}) is a loan, usually traced to Persian \emph{jāydād}, \textquotedblleft{}property, estate\textquotedblright{}.
\end{cousinweb}

\begin{grammarlens}[title={What you've built}]
One more word for this chapter.
\end{grammarlens}

\subsection*{Your turn}
\begin{itemize}
\raggedright
  \item \emph{Say it:} \emph{jāydād}
  \item \emph{Say it:} \emph{jāydād}, once more
  \item \emph{Say it:} say \emph{jāydād}
  \item \emph{Recall:} say the Marwadi for an owner, then the Marwadi for land, then say \emph{jāydād} again
\end{itemize}

\begin{tcolorbox}[breakable,colback=teal!4,colframe=teal!35!black,title={Before you move on}]
What does \mw{जायदाद} mean? (\textquotedblleft{}Property\textquotedblright{}.) Say it once more.
\end{tcolorbox}
\section[\texorpdfstring{\mw{पूंजी} (pū̃jī)}{पूंजी (pū̃jī)}]{\mw{पूंजी} (pū̃jī) — capital, savings}

\label{lesson:MW-C115-puji}



\begin{quote}
Before the new one: say the Marwadi for land, then the Marwadi for property.
\end{quote}

\subsection*{You'll want to know: \mw{पूंजी}}
\textbf{\mw{पूंजी}} — \emph{pū̃jī} — \textquotedblleft{}capital, savings\textquotedblright{}.

\begin{cousinweb}
\textbf{\mw{पूंजी}} (\emph{pū̃jī}) is a tadbhava, a word inherited from Sanskrit and reshaped by sound change. It comes from \emph{puñja}, \textquotedblleft{}heap, quantity\textquotedblright{}. CDIAL gives no Marwari form in this entry (CDIAL 8251); the match is through Hindi \emph{pū̃jī}. The sense moved from \textquotedblleft{}heap\textquotedblright{} to \textquotedblleft{}store of money\textquotedblright{}. CDIAL thinks the further history of the Sanskrit word doubtful.
\end{cousinweb}

\begin{grammarlens}[title={What you've built}]
One more word for this chapter.
\end{grammarlens}

\subsection*{Your turn}
\begin{itemize}
\raggedright
  \item \emph{Say it:} \emph{pū̃jī}
  \item \emph{Say it:} \emph{pū̃jī}, once more
  \item \emph{Say it:} say \emph{pū̃jī}
  \item \emph{Recall:} say the Marwadi for land, then the Marwadi for property, then say \emph{pū̃jī} again
\end{itemize}

\begin{tcolorbox}[breakable,colback=teal!4,colframe=teal!35!black,title={Before you move on}]
What does \mw{पूंजी} mean? (\textquotedblleft{}Capital, savings\textquotedblright{}.) Say it once more.
\end{tcolorbox}
\section[\texorpdfstring{\mw{कमाई} (kamāī)}{कमाई (kamāī)}]{\mw{कमाई} (kamāī) — earnings}

\label{lesson:MW-C115-kamai}



\begin{quote}
Before the new one: say the Marwadi for property, then the Marwadi for capital.
\end{quote}

\subsection*{You'll want to know: \mw{कमाई}}
\textbf{\mw{कमाई}} — \emph{kamāī} — \textquotedblleft{}earnings\textquotedblright{}.

\begin{cousinweb}
\textbf{\mw{कमाई}} (\emph{kamāī}) is a tadbhava, a word inherited and reshaped by sound change. It comes from \emph{*karmāpayati}, \textquotedblleft{}works, earns\textquotedblright{}, a form CDIAL reconstructs from its descendants. CDIAL gives no Marwari form in this entry (CDIAL 2897); the match is through Hindi \emph{kamānā} and Gujarati \emph{kamāvũ}. \textbf{\mw{कमाई}} (\emph{kamāī}) is the noun of the verb \textquotedblleft{}to earn\textquotedblright{}; CDIAL builds that verb on Sanskrit \emph{karman}, \textquotedblleft{}work\textquotedblright{}.
\end{cousinweb}

\begin{grammarlens}[title={What you've built}]
One more word for this chapter.
\end{grammarlens}

\subsection*{Your turn}
\begin{itemize}
\raggedright
  \item \emph{Say it:} \emph{kamāī}
  \item \emph{Say it:} \emph{kamāī}, once more
  \item \emph{Say it:} say \emph{kamāī}
  \item \emph{Recall:} say the Marwadi for property, then the Marwadi for capital, then say \emph{kamāī} again
\end{itemize}

\begin{tcolorbox}[breakable,colback=teal!4,colframe=teal!35!black,title={Before you move on}]
What does \mw{कमाई} mean? (\textquotedblleft{}Earnings\textquotedblright{}.) Say it once more.
\end{tcolorbox}
\section[\texorpdfstring{\mw{नौकरी} (naukrī)}{नौकरी (naukrī)}]{\mw{नौकरी} (naukrī) — a job}

\label{lesson:MW-C115-naukri}



\begin{quote}
Before the new one: say the Marwadi for capital, then the Marwadi for earnings.
\end{quote}

\subsection*{You'll want to know: \mw{नौकरी}}
\textbf{\mw{नौकरी}} — \emph{naukrī} — \textquotedblleft{}a job\textquotedblright{}.

\begin{cousinweb}
\textbf{\mw{नौकरी}} (\emph{naukrī}) is a loan, usually traced to Persian \emph{naukar}, \textquotedblleft{}servant\textquotedblright{}. Persian took \emph{naukar} from Mongolic \emph{nökör}, \textquotedblleft{}companion\textquotedblright{}; Marwadi added the ending \emph{-ī}, so \textquotedblleft{}the work of a servant\textquotedblright{}.
\end{cousinweb}

\begin{grammarlens}[title={What you've built}]
One more word for this chapter.
\end{grammarlens}

\subsection*{Your turn}
\begin{itemize}
\raggedright
  \item \emph{Say it:} \emph{naukrī}
  \item \emph{Say it:} \emph{naukrī}, once more
  \item \emph{Say it:} say \emph{naukrī}
  \item \emph{Recall:} say the Marwadi for capital, then the Marwadi for earnings, then say \emph{naukrī} again
\end{itemize}

\begin{tcolorbox}[breakable,colback=teal!4,colframe=teal!35!black,title={Before you move on}]
What does \mw{नौकरी} mean? (\textquotedblleft{}A job\textquotedblright{}.) Say it once more.
\end{tcolorbox}
